\begin{table}[tbp]
\centering
\small
\renewcommand{\arraystretch}{1.18}
\caption{Guida alle metriche e ai costi utilizzati nella valutazione.
Le frecce indicano il verso preferibile, a parità delle altre condizioni.
L'ECE dipende dal binning; Brier score e log-loss non misurano soltanto
la calibrazione.}
\label{tab:guida-metriche}
\begin{tabularx}{\textwidth}{@{}p{2.5cm}Yc@{}}
\toprule
\rowcolor{ThesisTableHeader}
\tablehead{Metrica} & \tablehead{Significato} & \tablehead{Verso} \\
\midrule
Accuracy & Quota di esempi classificati correttamente sul totale. & \(\uparrow\) \\
Precision & Per una classe: quota di assegnazioni a quella classe che sono corrette. & \(\uparrow\) \\
Recall & Per una classe: quota dei suoi esempi reali riconosciuta dal modello. & \(\uparrow\) \\
F1 per classe & Media armonica di precision e recall della classe. & \(\uparrow\) \\
Macro-F1 & Media delle F1 con uguale peso alle classi; metrica principale della tesi. & \(\uparrow\) \\
Balanced accuracy & Media dei recall; coincide con l'accuracy quando i supporti delle classi sono uguali. & \(\uparrow\) \\
Log-loss & Penalità sulle probabilità delle classi vere; cresce per errori molto confidenti. & \(\downarrow\) \\
Brier score & Errore quadratico delle probabilità rispetto alle etichette in forma one-hot. & \(\downarrow\) \\
ECE & Discrepanza aggregata tra confidenza e accuratezza, raggruppando le predizioni per fasce di confidenza. & \(\downarrow\) \\
Tempo di training & Somma dei tempi di epoca registrati; esclude la costruzione della cache di feature. & \(\downarrow\) \\
Tempo di test & Tempo di predizione sull'intero test dalle feature; non latenza end-to-end dalle immagini. & \(\downarrow\) \\
\bottomrule
\end{tabularx}
\end{table}

\begin{table}[tbp]
\centering
\small
\renewcommand{\arraystretch}{1.2}
\caption{Indicatori statistici: non sono metriche di classificazione.
Gli intervalli sono marginali, mentre Holm corregge i test all'interno
delle famiglie definite dal protocollo.}
\label{tab:guida-indicatori}
\begin{tabularx}{\textwidth}{@{}p{2.5cm}Y@{}}
\toprule
\rowcolor{ThesisTableHeader}
\tablehead{Indicatore} & \tablehead{Interpretazione} \\
\midrule
\(\Delta\) Q--B2 & Differenza media appaiata di Macro-F1. Positiva favorisce Q,
negativa B2; è espressa in punti percentuali. \\
IC95\% & Intervallo bootstrap per la quantità stimata, come media o differenza.
Non è l'intervallo che contiene il 95\% delle run. \\
\(p_{\mathrm{Holm}}\) & P-value corretto per i confronti multipli della famiglia.
Non è la probabilità che l'ipotesi nulla sia vera, né misura l'ampiezza dell'effetto. \\
\(d_z\) & Differenza media divisa per la deviazione standard delle differenze
appaiate: misura standardizzata dell'effetto. \\
\bottomrule
\end{tabularx}
\end{table}
